\documentclass[../../main.tex]{subfiles}
\begin{document}
\section{习题}\label{section:ch20-problems}

\begin{exercise}\label{exercise:ch20-01}
设 $G$ 是一个 Lie 群。
\begin{enumerate}[(a)]
\item 证明 $G$ 中单参数子群的像恰为维数不超过 $1$ 的连通 Lie 子群。
\item 证明每个单参数子群的像作为 Lie 群同构于下列之一：$\mathbb{R}$、$\mathbb{S}^1$ 或平凡群 $\{e\}$。
\end{enumerate}
\end{exercise}

\begin{exercise}\label{exercise:ch20-02}
计算交换 Lie 群 $\mathbb{R}^n$ 与 $\mathbb{T}^n$ 的指数映射。
\end{exercise}

\begin{exercise}\label{exercise:ch20-03}
设 $G$ 是一个 Lie 群，$A,B\subseteq\mathfrak g$ 是 $\operatorname{Lie}(G)$ 的互补线性子空间。证明由 $(X,Y)\mapsto(\exp X)(\exp Y)$ 给出的映射 $A\oplus B\to G$ 是从 $A\oplus B$ 中 $(0,0)$ 的某个邻域到 $G$ 中 $e$ 的某个邻域的微分同胚。（用于第 523 页。）
\end{exercise}

\begin{exercise}\label{exercise:ch20-04}
证明矩阵指数满足恒等式
\begin{align*}
\det e^A=e^{\operatorname{tr}A}.
\end{align*}
[提示：把命题\cref{proposition:ch20-指数性质}(g) 应用于 $\det\colon\operatorname{GL}(n,\mathbb{R})\to\mathbb{R}^*$。]
\end{exercise}

\begin{exercise}\label{exercise:ch20-05}
设 $a,b,c$ 为实数，$A,B,C$ 为 $\mathfrak{gl}(3,\mathbb{R})$ 中下列元素：
\begin{align*}
A=\begin{pmatrix}a&0&0\\0&b&0\\0&0&c\end{pmatrix},\qquad
B=\begin{pmatrix}0&a&b\\0&0&c\\0&0&0\end{pmatrix},\qquad
C=\begin{pmatrix}0&1&0\\-1&0&0\\0&0&0\end{pmatrix}.
\end{align*}
给出由 $A,B,C$ 生成的 $\operatorname{GL}(3,\mathbb{R})$ 单参数子群的显式公式（不是无穷级数）。
\end{exercise}

\begin{exercise}\label{exercise:ch20-06}
本习题表明连通 Lie 群的指数映射未必满射。
\begin{enumerate}[(a)]
\item 设 $A\in\operatorname{SL}(n,\mathbb{R})$ 对某个 $B\in\mathfrak{gl}(n,\mathbb{R})$ 形如 $e^B$。证明 $A$ 在 $\operatorname{SL}(n,\mathbb{R})$ 中有平方根，即存在矩阵 $C\in\operatorname{SL}(n,\mathbb{R})$ 使 $C^2=A$。
\item 令
\begin{align*}
A=\begin{pmatrix}-\tfrac12&0\\0&-2\end{pmatrix}.
\end{align*}
证明 $\exp\colon\mathfrak{sl}(2,\mathbb{R})\to\operatorname{SL}(2,\mathbb{R})$ 不满射，方法是表明 $A$ 不在其像中。[注：下一章习题 21-25 将证明 $\operatorname{SL}(2,\mathbb{R})$ 连通。]
\end{enumerate}
\end{exercise}

\begin{exercise}\label{exercise:ch20-07}
设 $G$ 是连通 Lie 群，$\mathfrak g$ 是它的 Lie 代数。
\begin{enumerate}[(a)]
\item 对任意 $X,Y\in\mathfrak g$，证明 $[X,Y]=0$ 当且仅当
\begin{align*}
(\exp tX)(\exp sY)=(\exp sY)(\exp tX)\qquad\text{对所有 }s,t\in\mathbb{R}.
\end{align*}
\item 证明 $G$ 交换当且仅当 $\mathfrak g$ 交换。
\item 给出当 $G$ 不连通时 (b) 的反例。
\end{enumerate}
\end{exercise}

\begin{exercise}\label{exercise:ch20-08}
设 $G$ 是一个 Lie 群。证明 $(\exp(X+Y))=(\exp X)(\exp Y)$ 对所有 $X,Y\in\operatorname{Lie}(G)$ 成立当且仅当 $G$ 的单位分支交换。[提示："若"方向，证明 $t\mapsto(\exp tX)(\exp tY)$ 是单参数子群；"仅当"方向，用习题 20-7。]
\end{exercise}

\begin{exercise}\label{exercise:ch20-09}
扩展命题\cref{proposition:ch20-乘积一阶} 的结果，证明在相同假设下存在光滑函数 $\widehat Z\colon(-\varepsilon,\varepsilon)\to\mathfrak g$，使
\begin{align*}
(\exp tX)(\exp tY)=\exp\bigl(t(X+Y)+\tfrac12t^2[X,Y]+t^3\widehat Z(t)\bigr).
\end{align*}
[注：有一个显式公式，称为 Baker--Campbell--Hausdorff 公式，给出满足 $(\exp tX)(\exp tY)=\exp\varphi(t)$ 的映射 $\varphi\colon(-\varepsilon,\varepsilon)\to\mathfrak g$ 的 Taylor 级数中所有项。上式给出该级数的前两项。完整公式见 [Var84]。]
\end{exercise}

\begin{exercise}\label{exercise:ch20-10}
设 $G$ 是一个 Lie 群，$S$ 是 $G$ 的 Lie 子群。证明 $S$ 的闭包也是 Lie 子群。由此推出 $G$ 的每个 Lie 子群或是 $G$ 的真嵌入子流形，或是某个真嵌入子流形的稠密子集。[注：这表明习题 7.20 中的子群 $S\subseteq\mathbb{T}^3$——一个真嵌入子群的稠密子群——是非嵌入 Lie 子群的典型。]
\end{exercise}

\begin{exercise}\label{exercise:ch20-11}
设 $G$ 与 $H$ 是 Lie 群。
\begin{enumerate}[(a)]
\item 证明每个连续同态 $\gamma\colon\mathbb{R}\to H$ 光滑。[提示：取 $V\subseteq\operatorname{Lie}(H)$ 为 $0$ 的邻域，使指数映射是从 $2V=\{2X\colon X\in V\}$ 到 $\exp(2V)$ 的微分同胚。取 $t_0$ 足够小，使 $|t|\leqslant t_0$ 时 $\gamma(t)\in\exp(V)$，令 $X_0\in V$ 使 $\gamma(t_0)=\exp X_0$。证明 $q$ 为二进有理数（即形如 $m/2^n$，$m,n\in\mathbb{Z}$）时 $\gamma(qt_0)=\exp(qX_0)$。]
\item 证明每个连续同态 $F\colon G\to H$ 光滑。[提示：证明存在映射 $\varphi\colon\operatorname{Lie}(G)\to\operatorname{Lie}(H)$ 使图表
\begin{align*}
\begin{CD}
\operatorname{Lie}(G) @>{\varphi}>> \operatorname{Lie}(H)\\
@V{\exp}VV @VV{\exp}V\\
G @>{F}>> H
\end{CD}
\end{align*}
交换；再用推论\cref{corollary:ch20-乘积极限} 证明 $\varphi$ 线性。]
\item 证明在 $G$ 上给定拓扑下，使 $G$ 成为 Lie 群的光滑结构只有一种。
\end{enumerate}
（用于第 556 页。）
\end{exercise}

\begin{exercise}\label{exercise:ch20-12}
设 $G$ 是一个 Lie 群。证明 $G$ 经由右平移在自身上作用的无穷小生成元是包含映射 $\operatorname{Lie}(G)\hookrightarrow\mathfrak X(G)$。
\end{exercise}

\begin{exercise}\label{exercise:ch20-13}
设 $\mathfrak g$ 是有限维 Lie 代数，$M$ 是光滑流形。若对每个 $p\in M$，形如 $\widehat X_p$（$\widehat X\in\widehat{(\,)}(\mathfrak g)$）的向量张成 $T_pM$，则称 Lie 代数作用 $\widehat{(\,)}\colon\mathfrak g\to\mathfrak X(M)$ 是\emph{可迁的}。证明 Lie 群 $G$ 在连通光滑流形 $M$ 上的光滑右作用可迁当且仅当其无穷小生成元可迁。[提示：证明若 Lie 代数作用可迁，则每个轨道开。]
\end{exercise}

\begin{exercise}\label{exercise:ch20-14}
设 $M$ 是光滑流形，$\mathfrak g\subseteq\mathfrak X(M)$ 是只由完全向量场组成的有限维 Lie 子代数。证明存在 Lie 群 $G$ 在 $M$ 上的光滑右作用，使 $\mathfrak g$ 为其无穷小生成元的像。对习题 8-20 中描述的 $\mathfrak X(\mathbb{R}^3)$ 的 Lie 子代数确定这样一个作用。
\end{exercise}

\begin{exercise}\label{exercise:ch20-15}
证明定理\cref{theorem:ch20-左作用基本定理}（左 Lie 代数作用基本定理）。[提示：用左作用与右作用间由 $g\cdot p=p\cdot g^{-1}$ 给出的一一对应。]
\end{exercise}

\begin{exercise}\label{exercise:ch20-16}
设 $G$ 是单连通 Lie 群，$\mathfrak g$ 是它的 Lie 代数。证明 $\mathfrak g$ 的每个表示都形如 $\rho_*\colon\mathfrak g\to\mathfrak{gl}(V)$，其中 $\rho\colon G\to\operatorname{GL}(V)$ 是 $G$ 的某个表示。
\end{exercise}

\begin{exercise}\label{exercise:ch20-17}
设 $G$ 是连通 Lie 群，$H$ 是任意 Lie 群，$\Phi,\Psi\colon G\to H$ 是 Lie 群同态且满足 $\Phi_*=\Psi_*\colon\operatorname{Lie}(G)\to\operatorname{Lie}(H)$。证明 $\Phi=\Psi$。（用于第 531 页。）
\end{exercise}

\begin{exercise}\label{exercise:ch20-18}
若 $\mathsf C$ 与 $\mathsf D$ 是范畴，协变函子 $\mathscr F\colon\mathsf C\to\mathsf D$ 称为\emph{范畴等价}，若 $\mathsf D$ 的每个对象都同构于某个 $\mathscr F(X)$（$X\in\operatorname{Ob}(\mathsf C)$），且映射 $\mathscr F\colon\operatorname{Hom}_{\mathsf C}(X,Y)\to\operatorname{Hom}_{\mathsf D}(\mathscr F(X),\mathscr F(Y))$ 对每对对象 $X,Y\in\operatorname{Ob}(\mathsf C)$ 双射。证明指派 $G\mapsto\operatorname{Lie}(G)$、$\varphi\mapsto\varphi_*$ 是单连通 Lie 群范畴 $\mathsf{SLie}$ 与有限维 Lie 代数范畴 $\mathsf{lie}$ 之间的范畴等价。
\end{exercise}

\begin{exercise}\label{exercise:ch20-19}
证明定理\cref{theorem:ch20-Lie基本定理}（Lie 的基本定理）。
\end{exercise}

\begin{exercise}\label{exercise:ch20-20}
设 $G$ 是连通 Lie 群，$\mathfrak g$ 是它的 Lie 代数。证明 $\operatorname{Ad}\colon G\to\operatorname{GL}(\mathfrak g)$ 的核是 $G$ 的\emph{中心}，即 $G$ 中与 $G$ 每个元素交换的元素集合。
\end{exercise}

\begin{exercise}\label{exercise:ch20-21}
证明 $\operatorname{GL}(n,\mathbb{R})$ 的伴随表示由 $\operatorname{Ad}(A)Y=AYA^{-1}$（$A\in\operatorname{GL}(n,\mathbb{R})$，$Y\in\mathfrak{gl}(n,\mathbb{R})$）给出。证明它不单。
\end{exercise}

\begin{exercise}\label{exercise:ch20-22}
若 $\mathfrak g$ 是 Lie 代数，$\mathfrak g$ 的\emph{中心}是所有满足 $[X,Y]=0$（对所有 $Y\in\mathfrak g$）的 $X\in\mathfrak g$ 的集合。设 $G$ 是连通 Lie 群。证明 $\operatorname{Lie}(G)$ 的中心是 $G$ 的中心的 Lie 代数。
\end{exercise}
\end{document}
